\subsection{REAL-VALUED FUNCTIONS}

DEFINITION 1. A mapping of a set X into the real line is called a real-valued function (or real function) defined on X.

By an abuse of language analogous to that mentioned in § 4, no. 2, mappings of X into \(\overline{\mathbf{R}}\) will also be called real-valued functions defined on X in this and the following section. Mappings of X into \(\mathbf{R}\) will be called finite real-valued functions.

If \(f\) and \(g\) are two real-valued functions defined on \(\mathbf{X}\) , the relation \(f \leqslant g\) is by definition equivalent to `` \(f(x) \leqslant g(x)\) for all \(x \in \mathbf{X}\) ;'' this relation is an ordering on the set \(\overline{\mathbf{R}}^{\mathbf{X}}\) of all real-valued functions on \(\mathbf{X}\) . The set \(\overline{\mathbf{R}}^{\mathbf{X}}\) ordered by this relation, is a lattice; for if \(f\) and \(g\) are any two real-valued functions, the function \(h\) defined by \(h(x) = \sup (f(x), g(x))\) for all \(x \in \mathbf{X}\) is the smallest of the real-valued functions on \(\mathbf{X}\) which are both \(\geqslant f\) and \(\geqslant g\) ; in accordance with general notation, we denote this function (which is the least upper bound of \(f\) and \(g\) in \(\overline{\mathbf{R}}^{\mathbf{X}}\) ) by \(\sup (f, g)\) . Similarly, the function whose value at each \(x \in \mathbf{X}\) is \(\inf (f(x), g(x))\) is denoted by \(\inf (f, g)\) .

Note that sup \((f, g)\) is the composition of the mapping

\[
(u, v) \mapsto \sup (u, v)
\]

of \(\overline{\mathbf{R}}\times \overline{\mathbf{R}}\) into \(\overline{\mathbf{R}}\) and the mapping \(x\to (f(x),g(x))\) of \(\mathbf{X}\) into \(\overline{\mathbf{R}}\times \overline{\mathbf{R}}\) . Similarly for inf \((f,g)\) .

A real-valued function \(f\) defined on a set \(\mathbf{X}\) is said to be bounded above (resp. bounded below) in \(\mathbf{X}\) , if \(f(\mathbf{X})\) is a subset of \(\mathbf{A}'' = [-\infty, +\infty[\) and is bounded above (resp. if \(f(\mathbf{X})\) is a subset of \(\mathbf{A}' = ] - \infty, +\infty]\) and is bounded below). \(f\) is said to be bounded in \(\mathbf{X}\) if it is bounded both above and below, that is if \(f(\mathbf{X})\) is a bounded subset of \(\mathbf{R}\) .

Every bounded function is therefore finite. The converse is false, as is shown by the function \(\mathbf{1} / x\) on \(\mathbb{R}_+^* = ]0, +\infty[\) .
% semantic-fragments: legacy-input-seam
\relax{}
